\section{Related literature}
\label{sec:literature}

The paper connects five literatures that are usually studied separately.

\paragraph{Endogenous innovation and economic growth.}
The first literature makes ideas production an economic activity. In
\citet{romer1990}, nonrival knowledge and imperfect competition sustain investment in
new varieties. \citet{aghionhowitt1992} model growth through quality-improving creative
destruction. \citet{jones1995} shows that diminishing returns in the production of ideas
remove scale effects and tie long-run growth to the growth of the research labor force.
Our starting point is closest to the semi-endogenous formulation: human researchers
are initially the scarce input that limits innovation. The departure is that research
labor can itself be produced with AI and compute.

\paragraph{Capital substitution and asymptotically linear growth.}
The capital-deepening benchmark belongs to a different growth tradition.
\citet{solow1956} and \citet{pitchford1960} show that a sufficiently high elasticity
of substitution can keep the marginal product of capital from vanishing in a model
with a constant saving rate. \citet{jonesmanuelli1990} characterize the broader class
of convex technologies that become asymptotically linear, and \citet{rebelo1991}
shows how an accumulable capital core can sustain endogenous growth without
increasing returns. \citet{duffypapageorgiou2000} emphasize the qualification that an
elasticity above one is necessary but not sufficient: asymptotic capital productivity
must also exceed the depreciation and population-growth requirement. Our benchmark
uses two capital services to make this comparison transparent. Its nested architecture
is close to \citet{kruselletral2000}, who combine structures capital in an outer
Cobb--Douglas layer with equipment capital and labor inputs inside a CES nest. Unlike
their quantitative model, our benchmark fixes the division of total capital and uses
it only to separate ordinary capital deepening from recursive improvement in AI
capability.

\paragraph{AI, automation, and the production of ideas.}
\citet{aghionjonesjones2019} examine AI as an automation technology and emphasize that
growth can remain constrained by tasks that are essential but difficult to automate.
\citet{acemoglu2025} also models AI at the task level and relates its aggregate
productivity effect to the share of tasks affected and the cost savings within those
tasks.
\citet{nordhaus2021} studies the substitutability conditions required for an economic
singularity. \citet{trammellkorinek2023} distinguish the automation of ordinary
production from the automation of R\&D and show that physical and technological
bottlenecks remain central even under transformative AI. \citet{jones2025rd} provides
a general framework in which AI's effect on research depends on the share of tasks it
can perform, its productivity in those tasks, and the strength of complementary
bottlenecks. Our model endogenizes the gradual composition of human and automated
research and embeds the ideas production function in a decentralized market
equilibrium.

\paragraph{Automation, wages, and the labor share.}
The effect of automation on labor income depends on both substitution and
productivity. \citet{acemoglurestrepo2019} distinguish the displacement effect of
automating existing tasks from the productivity and reinstatement effects that can
raise labor demand. In CES models, the unit-elasticity threshold determines whether
an increase in a factor reduces the income share of the other factor. This mechanism
underlies the account in \citet{karabarbounisneiman2014} of how cheaper investment
goods can reduce the labor share. A lower labor share, however, does not imply a lower
real wage. Using a nested CES and Korean firm and establishment data,
\citet{aumshin2024} estimate an elasticity of 0.6 between equipment and labor and an
elasticity of 1.6 between software and labor. The second estimate is a useful
empirical reference for our production elasticity $\sigma_{XL}$ because $X$ is a flow of
AI services rather than a stock of conventional equipment. It is not a direct
estimate of substitution between AI services and labor in the world economy.
\citet{autorkausik2026} establish more generally that automation can raise
wages while reducing labor's share in a competitive constant-returns economy.
\citet{yango2026} obtains, in a nested Cobb--Douglas--CES technology, the same two
thresholds that appear in our production block: the labor share falls when
$\sigma_{XL}>1$, while the wage rises when $\sigma_{XL}<1/\alpha$. Thus, for
$1<\sigma_{XL}<1/\alpha$, an expansion of AI simultaneously raises the wage and lowers
labor's share of income. \citet{trammellkorinek2023} place this ambiguity in the
broader transformative-AI literature: production automation can sharply reduce the
labor share, while the wage response continues to depend on technology and
accumulation.

We use the two-threshold result as a benchmark rather than claim it as a new
comparative static. Our departure is to make AI a flow of quality-adjusted services
supplied by an integrated monopolist and to make the underlying capability frontier
endogenous. Market power changes the equilibrium path of AI deployment, capital
accumulation, wages, and operating rents; frontier research determines how long the
economy remains in each production regime. This structure lets us study how service
regulation and research policy affect the wage and the labor share through different
static and dynamic channels.

\paragraph{AI services, inference compute, and quality adjustment.}
Using an existing model and producing a new one require separate measures of
inference and research compute. \citet{korinekmckelvey2026} combine inference compute
with intangible model capital and adjust the resulting output for quality.
\citet{demireretal2025} document prices per token within intelligence tiers, which
likewise implies that raw tokens are not homogeneous AI services. Our specification
$X_t=A_tU_t$ is an efficiency-units representation, not a technological identity:
$U_t$ measures the scale at which the current model is run, while $A_t$ converts
inference compute into quality-adjusted services. Task-based formulations such as
\citet{aghionjonesjones2019} and \citet{acemoglu2025} provide a richer alternative
when heterogeneity across uses of AI is central.

\paragraph{Market structure in AI.}
The industrial-organization literature emphasizes high fixed costs, scarce compute,
data, economies of scope, vertical integration, and feedback between model quality and
the surrounding ecosystem. \citet{korinekvipra2025} study how these forces can produce
concentration and tipping in foundation models. \citet{gans2024} emphasizes the role of
markets for training and input data in determining market power. Using model-usage
data, \citet{demireretal2025} document entry, falling prices, differentiation, and
turnover among leading models. The evidence is consistent with strong competition
today but does not rule out future concentration. Our extended model focuses on a
specific dynamic mechanism: a temporary lead becomes persistent when it lowers the
leader's cost of producing the next generation of frontier models.

Our baseline takes this concentration to its limiting case: one integrated frontier
developer supplies current services and conducts frontier research. Regulated access
prices are introduced as policy counterfactuals rather than as a second baseline
market structure. This formulation isolates the dynamic link between current
operating rents and the private incentive to improve the frontier.

\paragraph{Contribution relative to existing work.}
Several individual mechanisms in the model are established in the literature and
are not claimed as new. Semi-endogenous growth already ties long-run innovation to
population growth when human research labor is the expanding input. The possibility
of an economic singularity, the importance of essential tasks, and the distinction
between automating production and automating R\&D are central to
\citet{aghionjonesjones2019}, \citet{nordhaus2021}, and
\citet{trammellkorinek2023}. The two static thresholds separating the real wage from
labor's income share appear in recent automation models, including
\citet{autorkausik2026} and \citet{yango2026}. Nor is the general trade-off between
market power and innovation new. The paper uses these results as building blocks.

The first departure is to make the composition of AI research a continuous
equilibrium object inside a dynamic macroeconomic model. Relative to the task-based
research framework in \citet{jones2025rd}, the baseline suppresses heterogeneity
across research tasks but endogenizes the aggregate automation share through the
relative price of human and automated research. That share then enters the feedback
denominator governing capability growth. The task extension restores an extensive
automation margin while preserving the connection to the decentralized equilibrium.

The second departure is the formal comparison between human-essential and
AI-dominated research. With Cobb--Douglas production, the two feedback denominators
are
\[
    D_{\mathrm{CD}}
    =1-\phi-\eta\omega_M\frac{\omega_X}{\omega_L},
    \qquad
    D_{\mathrm{AI}}
    =1-\phi-\eta\frac{\omega_X}{\omega_L}.
\]
Their difference identifies a parameter region in which the human-essential
constant-growth candidate has a positive finite rate but the AI-dominated candidate
does not. In the reduced proportional-allocation system, the changing CES research
share also yields an exact path condition for finite-time singularity. That condition
characterizes the reduced system; it is not an existence theorem for the fully
optimizing equilibrium.
The fully Cobb--Douglas counterfactual also connects the paper to the recursive-growth
mechanism in \citet{aghionjonesjones2019}. Within the reduced system, at unit
substitution elasticities in both nests, the sign of $D_{\mathrm{CD}}$ alone
separates a stable positive growth rate,
double-exponential capability, and a finite-time singularity. Gross substitution in
final production is therefore one route to acceleration, not a necessary condition
for it.

The third departure concerns the interaction between recursive innovation and the
production CES. The paper proves that, for every fixed $\sigma_{XL}>1$, an interior
steady state with positive constant capability growth cannot exist: such a path would
make capability and the return to capital diverge, contradicting the household Euler
equation. It also shows that the unit-elasticity boundary is
a singular perturbation: finite allocations are continuous in $\sigma_{XL}$, while
the infinite-horizon limits do not commute and the scale required for AI dominance
diverges as the elasticity approaches one. This result complements the broader phase
diagrams in \citet{nordhaus2021} and \citet{trammellkorinek2023} by deriving the
transition in an economy with quality-adjusted AI services, capital accumulation,
and monopoly pricing.

The fourth departure is to connect research composition, research employment, and
market structure. The ratio $H/M$ can fall while $H/N$ rises because automation
changes the human intensity of research and the scale of the research industry at
the same time. The integrated developer links those employment dynamics to current
AI-service rents and frontier investment. This combination differs from the
market-structure analyses in \citet{gans2024} and \citet{korinekvipra2025}, which explain
the sources of concentration but do not embed the service provider's research choice
in the same aggregate transition system. We therefore describe the contribution as
the joint equilibrium characterization of these mechanisms rather than claim
priority for each ingredient separately.
